% !TEX program = xelatex
% 공통수학2 · Ⅱ. 집합과 명제 · 1. 집합 · 01 집합 · 1/2차시
% 학생용 활동지 (답안 없음)
\documentclass[11pt,a4paper]{article}
\usepackage{kotex}
\usepackage{iftex}
\ifPDFTeX\else
  \setmainhangulfont{Noto Serif CJK KR}
  \setsanshangulfont{Noto Sans CJK KR}
\fi
\usepackage[margin=2.1cm]{geometry}
\usepackage{parskip}
\usepackage{enumitem}
\usepackage{array}
\usepackage{tabularx}
\usepackage{xcolor}
\usepackage{amssymb}
\usepackage{tikz}
\usepackage[most]{tcolorbox}
\usepackage{fancyhdr}
\usepackage{titlesec}

\definecolor{goldc}{HTML}{B07C22}
\definecolor{goldstrong}{HTML}{8A5F16}
\definecolor{indigoc}{HTML}{2B3B57}
\definecolor{indigosoft}{HTML}{6E82A6}
\definecolor{linec}{HTML}{D6DACB}
\definecolor{surface2}{HTML}{E4E8DC}
\definecolor{writeline}{HTML}{C7CCBA}
\definecolor{inksoft}{HTML}{52604F}

\titleformat{\section}{\normalfont\sffamily\bfseries\large\color{indigoc}}{\thesection}{0.6em}{}
\titlespacing{\section}{0pt}{14pt}{6pt}
\newtcolorbox{goalbox}{enhanced, breakable, colback=white, colframe=goldc, boxrule=0.5pt, leftrule=3pt, arc=2pt, left=10pt, right=10pt, top=8pt, bottom=8pt}
\newtcolorbox{sheetbox}{enhanced, breakable, colback=white, colframe=linec, boxrule=0.6pt, arc=3pt, left=11pt, right=11pt, top=9pt, bottom=9pt}
\newcommand{\writespace}[1]{%
  \par\nobreak\vspace{3pt}
  \foreach \n in {1,...,#1} {%
    \noindent\textcolor{writeline}{\rule{\linewidth}{0.4pt}}\par\vspace{0.62cm}%
  }
}
\newcommand{\fillblank}[1]{\makebox[#1]{\hrulefill}}
\pagestyle{fancy}\fancyhf{}\renewcommand{\headrulewidth}{0pt}
\fancyfoot[L]{\footnotesize\color{inksoft} 공통수학2 · 2026학년도 2학기 · 지평선고등학교}
\fancyfoot[R]{\footnotesize\color{inksoft} 다음 시간 --- 01 집합(2) 원소의 개수}

\begin{document}

\noindent
\begin{minipage}[b]{0.62\linewidth}
  {\sffamily\footnotesize\color{indigosoft} 공통수학2 · Ⅱ. 집합과 명제 · 1. 집합}\\[2pt]
  {\sffamily\bfseries\Large 01 집합 \textperiodcentered\ 18차시}
\end{minipage}\hfill
\begin{minipage}[b]{0.36\linewidth}
  \raggedleft\small\color{inksoft}
  반\ \fillblank{1.6cm} \quad 번호\ \fillblank{1.6cm}\\[4pt]
  이름\ \fillblank{3.6cm}
\end{minipage}

\vspace{4pt}
\noindent{\color{goldc}\rule{\linewidth}{2.2pt}}
\vspace{10pt}

\begin{goalbox}
\textbf{오늘의 목표}\\
집합과 원소의 뜻을 이해하고, 원소나열법·조건제시법·벤 다이어그램으로 집합을 표현할 수 있다.
\vspace{4pt}
\small\color{inksoft}
$\square$ 자신 있음 \quad $\square$ 조금 헷갈림 \quad $\square$ 처음 봄 \hfill (수업 전 나의 상태에 표시)
\end{goalbox}

\section{생각 열기 --- 기준이 분명한 모임}

\begin{sheetbox}
우리나라 멸종 위기 야생 생물 Ⅰ급: 늑대, 수달, 혹고니, 수원청개구리, 감돌고기

\begin{enumerate}[leftmargin=1.4em, itemsep=2pt]
  \item 포유류에 속하는 동물을 모두 말해 보자.
  \item 귀여운 동물을 모두 말해 보고, 짝의 답과 비교해 보자. 왜 답이 다를 수 있을까?
\end{enumerate}
\writespace{3}
\end{sheetbox}

\section{개념 정리 --- 집합과 원소}

\begin{sheetbox}
어떤 기준에 따라 대상을 \fillblank{4cm} 정할 수 있을 때, 그 대상들의 모임을 \textbf{집합}이라 하고, 집합을 이루는 대상 하나하나를 그 집합의 \textbf{원소}라고 한다.

\vspace{4pt}
$a$가 집합 $A$의 원소일 때 $a\ \fillblank{1cm}\ A$, 원소가 아닐 때 $a\ \fillblank{1cm}\ A$로 나타낸다.

\vspace{8pt}
\textbf{집합의 표현}
\begin{itemize}[leftmargin=1.6em, itemsep=3pt]
  \item 원소나열법: 예) 8의 약수의 집합 $=\{\ \fillblank{4cm}\ \}$
  \item 조건제시법: 예) 위 집합 $=\{x\mid\ \fillblank{4cm}\ \}$
  \item 벤 다이어그램: 그림으로 나타내기
\end{itemize}
\end{sheetbox}

\section{연습 문제}

\begin{sheetbox}
\textbf{1.} 다음에서 집합인 것을 모두 찾고, 그 집합의 원소를 구하시오.\\
\hspace*{1.6em}⑴ 9의 약수의 모임 \quad ⑵ 인기 있는 수영 선수의 모임 \quad ⑶ 높은 빌딩의 모임 \quad ⑷ 방정식 $4x-5=0$의 해의 모임
\writespace{3}

\vspace{6pt}
\textbf{2.} 실수 전체의 집합을 $R$이라 할 때, 다음 $\square$ 안에 $\in$, $\notin$ 중 알맞은 것을 써넣으시오.\\
\hspace*{1.6em}⑴ $0\ \square\ R$ \qquad ⑵ $\sqrt[3]{21}\ \square\ R$ \qquad ⑶ $-i\ \square\ R$
\writespace{2}

\vspace{6pt}
\textbf{3.} 다음 집합에서 원소나열법으로 나타낸 것은 조건제시법으로, 조건제시법으로 나타낸 것은 원소나열법으로 나타내시오.\\
\hspace*{1.6em}⑴ $\{1,2,3,4,6,12\}$ \qquad ⑵ $\{5,10,15,\dots\}$\\
\hspace*{1.6em}⑶ $\{x\mid(x-3)(x-5)=0\}$ \qquad ⑷ $\{x\mid x$는 100 이하의 홀수$\}$
\writespace{4}
\end{sheetbox}

\section{사고의 가시화 --- See · Think · Wonder}

\noindent{\footnotesize\color{inksoft} 오늘 배운 `집합'을 떠올리며 아래 세 칸을 채워 보자.}\\[6pt]
\begin{tabularx}{\linewidth}{@{}XXX@{}}
\textbf{\footnotesize\color{indigoc} See --- 무엇을 보았나?} &
\textbf{\footnotesize\color{indigoc} Think --- 무엇을 생각했나?} &
\textbf{\footnotesize\color{indigoc} Wonder --- 무엇이 궁금한가?} \\[4pt]
\end{tabularx}
\noindent
\begin{minipage}[t]{0.32\linewidth}\writespace{3}\end{minipage}\hfill
\begin{minipage}[t]{0.32\linewidth}\writespace{3}\end{minipage}\hfill
\begin{minipage}[t]{0.32\linewidth}\writespace{3}\end{minipage}

\section{마무리 성찰}

\begin{sheetbox}
\begin{minipage}[t]{0.48\linewidth}
{\footnotesize\color{inksoft} 오늘 배운 것을 한 문장으로}
\writespace{2}
\end{minipage}\hfill
\begin{minipage}[t]{0.48\linewidth}
{\footnotesize\color{inksoft} 감사 일기 / 유념 한 줄}
\writespace{2}
\end{minipage}
\end{sheetbox}

\end{document}
